\begin{table}[!htbp]
\centering
\caption{Main EB school-value IV estimates controlling for SAE priority status}
\label{tab:scalar-school-value-iv-main-seven-eb-prioritario}
\resizebox{\textwidth}{!}{%
\begin{tabular}{lccccccc}
\toprule
 & \multicolumn{3}{c}{Exams} & \multicolumn{3}{c}{Higher ed. choices} & Financial aid \\
\cmidrule(lr){2-4} \cmidrule(lr){5-7} \cmidrule(lr){8-8}
 & Math & Verbal & Exam taking & High-premium inst. & High-premium field & Program income & Fin. aid app. \\
\midrule
$\psi^{EB}$ & 1.185*** & 1.016*** & 0.480*** & 0.660*** & 0.861*** & 1.057*** & 0.584*** \\
 & (0.097) & (0.098) & (0.052) & (0.120) & (0.087) & (0.131) & (0.098) \\
N & 92,361 & 93,980 & 135,844 & 94,487 & 94,487 & 94,487 & 135,844 \\
\bottomrule
\end{tabular}
}
\par\medskip
\footnotesize
\begin{minipage}{\textwidth}
Notes: This robustness specification adds the SAE low-income priority indicator, \textit{prioritario}, to the controls used in the corresponding main specification. It also controls for the DA-probability expected value of the same EB measure, cohort, grade-4 SIMCE math and language, gender, and age. The sample is restricted to students with observed priority status. Heteroskedasticity-robust standard errors are reported in parentheses. $^{*}p<0.10$, $^{**}p<0.05$, $^{***}p<0.01$.
\end{minipage}
\end{table}
